% GENERATED FILE. Edit canonical lessons, then run npm run generate:books.
% canonical-source-hash: fnv1a64:6469d791c9316478
% canonical-lessons: SA-C144-bhajati, SA-C144-tolayati, SA-C144-utpatati, SA-C144-niksipati, SA-C144-yojayati

\chapter{Shares, Weighs, Flies up, Puts in, Joins}
\label{ch:sa-shares-weighs-flies-up-puts-in-joins}
\hlchaptermodality{144}

\begin{chapteropening}
\textbf{By the end of this chapter:} \emph{I can say shares, weighs, flies up, puts in and joins.}

Complete the last lesson of chapter 144: I can say shares, weighs, flies up, puts in and joins.
\end{chapteropening}

\section[\texorpdfstring{\sk{भजति} (bhajati)}{भजति (bhajati)}]{\sk{भजति} (bhajati) — shares, worships}

\label{lesson:SA-C144-bhajati}



\begin{quote}
Before the new one: say the Sanskrit for is afraid, then the Sanskrit for meditates.
\end{quote}

\subsection*{You'll want to know: \sk{भजति}}
\textbf{\sk{भजति}} — \emph{bhajati} — \textquotedblleft{}shares, worships\textquotedblright{}.

\begin{grammarlens}[title={What you've built}]
One more word for this chapter.
\end{grammarlens}

\subsection*{Your turn}
\begin{itemize}
\raggedright
  \item \emph{Say it:} \emph{bhajati}
  \item \emph{Say it:} \emph{bhajati}, once more
  \item \emph{Say it:} say \emph{bhajati}
  \item \emph{Recall:} say the Sanskrit for is afraid, then the Sanskrit for meditates, then say \emph{bhajati} again
\end{itemize}

\begin{tcolorbox}[breakable,colback=teal!4,colframe=teal!35!black,title={Before you move on}]
What does \sk{भजति} mean? (\textquotedblleft{}Shares, worships\textquotedblright{}.) Say it once more.
\end{tcolorbox}
\section[\texorpdfstring{\sk{तोलयति} (tolayati)}{तोलयति (tolayati)}]{\sk{तोलयति} (tolayati) — weighs}

\label{lesson:SA-C144-tolayati}



\begin{quote}
Before the new one: say the Sanskrit for meditates, then the Sanskrit for shares.
\end{quote}

\subsection*{You'll want to know: \sk{तोलयति}}
\textbf{\sk{तोलयति}} — \emph{tolayati} — \textquotedblleft{}weighs\textquotedblright{}.

\begin{grammarlens}[title={What you've built}]
One more word for this chapter.
\end{grammarlens}

\subsection*{Your turn}
\begin{itemize}
\raggedright
  \item \emph{Say it:} \emph{tolayati}
  \item \emph{Say it:} \emph{tolayati}, once more
  \item \emph{Say it:} say \emph{tolayati}
  \item \emph{Recall:} say the Sanskrit for meditates, then the Sanskrit for shares, then say \emph{tolayati} again
\end{itemize}

\begin{tcolorbox}[breakable,colback=teal!4,colframe=teal!35!black,title={Before you move on}]
What does \sk{तोलयति} mean? (\textquotedblleft{}Weighs\textquotedblright{}.) Say it once more.
\end{tcolorbox}
\section[\texorpdfstring{\sk{उत्पतति} (utpatati)}{उत्पतति (utpatati)}]{\sk{उत्पतति} (utpatati) — flies up, jumps up}

\label{lesson:SA-C144-utpatati}



\begin{quote}
Before the new one: say the Sanskrit for shares, then the Sanskrit for weighs.
\end{quote}

\subsection*{You'll want to know: \sk{उत्पतति}}
\textbf{\sk{उत्पतति}} — \emph{utpatati} — \textquotedblleft{}flies up, jumps up\textquotedblright{}.

\begin{grammarlens}[title={What you've built}]
One more word for this chapter.
\end{grammarlens}

\subsection*{Your turn}
\begin{itemize}
\raggedright
  \item \emph{Say it:} \emph{utpatati}
  \item \emph{Say it:} \emph{utpatati}, once more
  \item \emph{Say it:} say \emph{utpatati}
  \item \emph{Recall:} say the Sanskrit for shares, then the Sanskrit for weighs, then say \emph{utpatati} again
\end{itemize}

\begin{tcolorbox}[breakable,colback=teal!4,colframe=teal!35!black,title={Before you move on}]
What does \sk{उत्पतति} mean? (\textquotedblleft{}Flies up, jumps up\textquotedblright{}.) Say it once more.
\end{tcolorbox}
\section[\texorpdfstring{\sk{निक्षिपति} (nikṣipati)}{निक्षिपति (nikṣipati)}]{\sk{निक्षिपति} (nikṣipati) — puts in, deposits}

\label{lesson:SA-C144-niksipati}



\begin{quote}
Before the new one: say the Sanskrit for weighs, then the Sanskrit for flies up.
\end{quote}

\subsection*{You'll want to know: \sk{निक्षिपति}}
\textbf{\sk{निक्षिपति}} — \emph{nikṣipati} — \textquotedblleft{}puts in, deposits\textquotedblright{}.

\begin{grammarlens}[title={What you've built}]
One more word for this chapter.
\end{grammarlens}

\subsection*{Your turn}
\begin{itemize}
\raggedright
  \item \emph{Say it:} \emph{nikṣipati}
  \item \emph{Say it:} \emph{nikṣipati}, once more
  \item \emph{Say it:} say \emph{nikṣipati}
  \item \emph{Recall:} say the Sanskrit for weighs, then the Sanskrit for flies up, then say \emph{nikṣipati} again
\end{itemize}

\begin{tcolorbox}[breakable,colback=teal!4,colframe=teal!35!black,title={Before you move on}]
What does \sk{निक्षिपति} mean? (\textquotedblleft{}Puts in, deposits\textquotedblright{}.) Say it once more.
\end{tcolorbox}
\section[\texorpdfstring{\sk{योजयति} (yojayati)}{योजयति (yojayati)}]{\sk{योजयति} (yojayati) — joins, arranges}

\label{lesson:SA-C144-yojayati}



\begin{quote}
Before the new one: say the Sanskrit for flies up, then the Sanskrit for puts in.
\end{quote}

\subsection*{You'll want to know: \sk{योजयति}}
\textbf{\sk{योजयति}} — \emph{yojayati} — \textquotedblleft{}joins, arranges\textquotedblright{}.

\begin{grammarlens}[title={What you've built}]
One more word for this chapter.
\end{grammarlens}

\subsection*{Your turn}
\begin{itemize}
\raggedright
  \item \emph{Say it:} \emph{yojayati}
  \item \emph{Say it:} \emph{yojayati}, once more
  \item \emph{Say it:} say \emph{yojayati}
  \item \emph{Recall:} say the Sanskrit for flies up, then the Sanskrit for puts in, then say \emph{yojayati} again
\end{itemize}

\begin{tcolorbox}[breakable,colback=teal!4,colframe=teal!35!black,title={Before you move on}]
What does \sk{योजयति} mean? (\textquotedblleft{}Joins, arranges\textquotedblright{}.) Say it once more.
\end{tcolorbox}
